\documentclass{../../../unipd-notes}

\unipdsetup{
  course = {Fondamenti di Ingegneria Informatica},
  short-course = {Fondamenti di Ingegneria},
  document-type = {Appunti delle lezioni}
}

\begin{document}
\chapter{Notazione matematica}

Una formula in linea, come $f(x)=x^2$, mantiene la stessa tipografia delle formule isolate.

\section{Equazioni e matrici}
\begin{equation}\label{eq:quadratic}
  x_{1,2}=\frac{-b\pm\sqrt{b^2-4ac}}{2a}
\end{equation}
L'\cref{eq:quadratic} è numerata automaticamente. Un calcolo allineato usa
\begin{align}
  (x+1)^2 &= x^2+2x+1,\\
  (x-1)^2 &= x^2-2x+1.
\end{align}

Siano
\[
  \vect{x}=\begin{bmatrix}x_1\\x_2\end{bmatrix},\qquad
  \mat{A}=\begin{bmatrix}1&2\\0&1\end{bmatrix},\qquad
  \mat{A}\transpose\vect{x}=\vect{b}.
\]

\section{Insiemi, probabilità e unità}
La funzione $\functionmap{g}{\R^2}{\R}$ è definita su
$\set{\vect{x}\in\R^2 \suchthat \norm*{\vect{x}}\leq 1}$.
In probabilità si scrive
\[
  \Prob\left(A\given B\right)=\frac{\Prob(A\cap B)}{\Prob(B)},
  \qquad \Expected[X],\quad \Var(X),\quad \Cov(X,Y).
\]
Le inclusioni $\N\subset\Z\subset\Q\subset\R\subset\C$ riassumono gli insiemi numerici principali. Per $z=3+4\imaginaryunit$ valgono $\abs*{z}=5$ e $\exponential^{\imaginaryunit\pi}=-1$, mentre $\inner*{\vect{x},\vect{y}}$ indica il prodotto interno. Inoltre,
\[
  \int_0^1 x^2\,\diff x=\frac{1}{3},\qquad
  \floor*{2.7}=2,\qquad \ceil*{2.1}=3.
\]
Per una matrice quadrata si usano $\diag(1,2,3)$ e $\rank(\mat{A})$; un problema di ottimizzazione può essere scritto come
\[
  x_{\min}=\argmin_{x\in\R} f(x),\qquad
  x_{\max}=\argmax_{x\in\R} f(x).
\]
Le unità seguono una forma uniforme: $g=\qty{9.81}{\metre\per\second\squared}$ e $R=\qty{4.7}{\kilo\ohm}$.
\end{document}
